% =====================================================================
% Cross-Version Adapter Portability via Representational Alignment
% ICLR 2027 submission skeleton
% Compile: latexmk -pdf main.tex   (uses iclr2027 style if present,
%          falls back to article so the skeleton always builds)
% =====================================================================
\documentclass[11pt]{article}

% --- Try ICLR style if dropped in; otherwise stay self-contained ------
\IfFileExists{iclr2027_conference.sty}{%
  \usepackage{iclr2027_conference}
}{%
  \usepackage[margin=1in]{geometry}
}

\usepackage{amsmath,amssymb}
\usepackage{graphicx}
\usepackage{booktabs}
\usepackage{natbib}
\usepackage{hyperref}
\usepackage{xcolor}

\title{Cross-Version Adapter Portability via Representational Alignment}
\author{Anonymous}
\date{}

\begin{document}
\maketitle

\begin{abstract}
LoRA adapters are entangled with the representation space of the base model
they were trained on, so a new base-model version forces a full retrain on the
original task data. We formalize \emph{cross-version adapter portability}:
transplanting a LoRA adapter from base $M_v$ onto a same-family successor
$M_{v+1}$ using only the two base checkpoints, the adapter, and a small
\emph{unlabeled} anchor set --- no task data, no retraining. We frame this as a
representational alignment problem, empirically characterize the geometric drift
between base versions with CKA, and derive a closed-form per-module linear
transplant $\Delta W' = P^{\mathrm{out}}\Delta W\,P^{\mathrm{in}}$ that preserves
the adapter's low-rank structure and adds no inference cost. Across
%TODO: fill final numbers
$k$ fine-tuning tasks and a family version pair, alignment recovers a substantial
fraction of the gap between naive transfer and a retrained oracle at a small
fraction of oracle compute.
\end{abstract}

% =====================================================================
% Introduction + Contribution statement
% =====================================================================
\section{Introduction}
\label{sec:intro}

Parameter-efficient fine-tuning with LoRA~\citep{hu2022lora} has made it cheap
to specialize a base model for a task: practitioners maintain libraries of
adapters layered over a frozen foundation checkpoint. But base models are not
static. When a vendor ships a new version of a family --- Gemma~3 to Gemma~4,
Qwen2 to Qwen2.5 --- every adapter trained on the old base must, in current
practice, be retrained from scratch on the original task data. That data is
often private, expensive to re-collect, or simply gone, and the retraining cost
scales with the number of adapters an organization maintains.

We ask whether a fine-tuning delta can instead be treated as a \emph{portable
object}: transplanted onto a new base version without retraining on task data.
The obstacle is that a LoRA adapter is entangled with the base model's
representation space (\S\ref{sec:problem}) --- it reads and writes in the old
model's activation geometry, which the new version does not preserve. Existing
weight-space transfer methods do not address this: task arithmetic and task
vectors~\citep{ilharco2023taskvectors}, TIES-Merging~\citep{yadav2023ties}, and
DARE~\citep{yu2024dare} all assume a \emph{fixed, shared} representation space
across the weights being combined, and break under version drift.

Our framing is that cross-version portability is a \emph{representational
alignment} problem. We measure the drift between base versions, learn a
lightweight linear map between their activation spaces from a small unlabeled
anchor set, and use it to re-express the adapter in the new model's geometry
while preserving its low-rank structure.

\paragraph{Contributions.}
\begin{itemize}
  \item \textbf{A novel problem: cross-version adapter portability.} We
        formalize transplanting a LoRA adapter from base $M_v$ onto a
        same-family successor $M_{v+1}$ using only the two base checkpoints, the
        adapter, and a small unlabeled anchor set --- no task data and no
        retraining (\S\ref{sec:problem}). To our knowledge this problem is
        unaddressed by prior model-merging or task-arithmetic work.
  \item \textbf{An empirical study of representational drift across model
        versions.} Using CKA over anchor activations, we characterize how, where,
        and how much representation geometry shifts between successive base
        versions, and show this drift predicts the failure of naive adapter
        transfer (\S\ref{sec:drift}).
  \item \textbf{A simple, closed-form alignment method.} We derive a per-module
        linear transplant $\Delta W' = P^{\mathrm{out}} \Delta W\, P^{\mathrm{in}}$
        (ridge / orthogonal-Procrustes) that preserves rank $r$ and thus adds no
        inference cost, fit from a single anchor-set forward pass through each
        base model (\S\ref{sec:method}).
  \item \textbf{Evaluation across tasks and a version pair.} We evaluate on
        2--3 fine-tuning tasks and 1--2 family version pairs against a naive-transfer
        lower bound and a retrained oracle upper bound, reporting the fraction of
        the oracle--naive gap recovered at a fraction of oracle compute
        (\S\ref{sec:experiments}).
  \item \textbf{Open release.} We release code, anchor sets, and the alignment
        maps to support reproducibility and follow-on work (\S\ref{sec:repro}).
\end{itemize}


\section{Related Work}
\label{sec:related}
% Task arithmetic / task vectors (Ilharco et al. 2023); TIES-Merging (Yadav et al.
% 2023); DARE (Yu et al. 2024) -- all assume shared/fixed representation space.
% Model stitching (Lenc & Vedaldi 2015; Bansal et al. 2021) -- learns alignment
% between layer reps but for stitching, not weight-level adapter transfer.
% Representation similarity: CKA (Kornblith et al. 2019); Platonic Representation
% Hypothesis (Huh et al. 2024). LoRA/PEFT (Hu et al. 2022). Knowledge distillation
% (requires inference + task signal). Position: none address cross-version
% weight-level adapter transfer without task data.

% =====================================================================
% Problem Formulation
% Cross-Version Adapter Portability via Representational Alignment
% =====================================================================
\section{Problem Formulation}
\label{sec:problem}

\subsection{Setup and Notation}

Let $\mathcal{F}$ denote a model \emph{family} (e.g.\ Gemma, Qwen) and let
$M_v$ and $M_{v+1}$ be two \emph{base} (pre-trained, instruction-tuned, or
foundation) checkpoints from $\mathcal{F}$ released at successive versions
$v$ and $v{+}1$. The two checkpoints typically share architecture and
tokenizer but differ in weights due to retraining; we make no assumption
that they share a representation space.

A \emph{LoRA adapter}~\citep{hu2022lora} $A$ trained on top of $M_v$ is a set
of low-rank updates over a collection of adapted linear modules
$\ell \in \mathcal{L}$. Module $\ell$ realizes a linear map with frozen
weight $W_\ell \in \mathbb{R}^{d^{\mathrm{out}}_\ell \times d^{\mathrm{in}}_\ell}$,
and LoRA augments it with
\begin{equation}
  \Delta W_\ell \;=\; \tfrac{\alpha}{r}\, B_\ell A_\ell,
  \qquad
  B_\ell \in \mathbb{R}^{d^{\mathrm{out}}_\ell \times r},\;
  A_\ell \in \mathbb{R}^{r \times d^{\mathrm{in}}_\ell},
\label{eq:lora}
\end{equation}
with rank $r \ll \min(d^{\mathrm{in}}_\ell, d^{\mathrm{out}}_\ell)$ and scale
$\alpha$. The fine-tuned model is denoted $M_v \oplus A$, i.e.\ each adapted
module computes $y = (W_\ell + \Delta W_\ell)\,x$.

The adapter $A$ is trained on a task dataset $\mathcal{D}_{\mathrm{task}}$
(potentially private, expensive, or no longer available). We additionally
assume access to a small, \emph{unlabeled} \emph{anchor} set
$\mathcal{D}_{\mathrm{anchor}}$ ($|\mathcal{D}_{\mathrm{anchor}}| \ll
|\mathcal{D}_{\mathrm{task}}|$, e.g.\ a few thousand generic text sequences)
that can be passed through both base checkpoints.

\subsection{Objective}

Let $\mathcal{P}(M \oplus A)$ be task performance of model $M$ equipped with
adapter $A$. We seek a transformation $T$ that transplants $A$ from $M_v$ onto
$M_{v+1}$ \emph{without access to $\mathcal{D}_{\mathrm{task}}$}:
\begin{equation}
  T^\star \;=\;
  \arg\max_{T \in \mathcal{T}}\;
  \mathcal{P}\!\left(M_{v+1} \oplus T(A)\right),
  \qquad
  \text{using only } M_v,\, M_{v+1},\, A,\, \mathcal{D}_{\mathrm{anchor}}.
\label{eq:objective}
\end{equation}
We compare against two reference points: the \emph{naive transfer} baseline
$T = \mathrm{id}$ (drop $A$ onto $M_{v+1}$ unchanged), and the \emph{oracle}
$A^\star = \arg\max_A \mathcal{P}(M_{v+1}\oplus A)$ obtained by retraining on
$\mathcal{D}_{\mathrm{task}}$. A successful $T$ closes the gap between naive
transfer and the oracle at a small fraction of the oracle's compute, and
without task labels.

\subsection{Why Naive Transfer Fails: Representational Drift}

Equation~\eqref{eq:lora} shows the adapter acts on the module's \emph{input}
activation and writes into its \emph{output} activation. Concretely, for input
$x$ in the input representation space of $M_v$, the adapter contributes
$\tfrac{\alpha}{r} B_\ell (A_\ell x)$. The adapter is therefore \emph{bound}
to the geometry of $M_v$'s activation spaces: $A_\ell$ reads coordinates in
$M_v$'s input basis and $B_\ell$ writes coordinates in $M_v$'s output basis.

When the same module's activation geometry differs in $M_{v+1}$, the directions
$A_\ell$ reads from and $B_\ell$ writes to no longer correspond to the same
semantic features, so the update is applied in the wrong basis. We quantify
this \emph{representational drift} per module with linear Centered Kernel
Alignment~\citep{kornblith2019cka}: collecting input activation matrices
$H^{v}_\ell, H^{v+1}_\ell \in \mathbb{R}^{N \times d^{\mathrm{in}}_\ell}$ over
$\mathcal{D}_{\mathrm{anchor}}$ ($N$ tokens),
\begin{equation}
  \mathrm{drift}_\ell \;=\; 1 - \mathrm{CKA}\!\left(H^{v}_\ell, H^{v+1}_\ell\right),
  \qquad
  \mathrm{CKA}(X,Y) = \frac{\lVert Y^\top X \rVert_F^2}
       {\lVert X^\top X \rVert_F \,\lVert Y^\top Y \rVert_F}.
\label{eq:cka}
\end{equation}
A central empirical claim of this paper (\S\ref{sec:drift}) is that
$\mathrm{drift}_\ell > 0$ across versions and grows with depth, and that this
drift predicts the failure of naive transfer.

\subsection{Representational Alignment}

We restrict the transformation class $\mathcal{T}$ to act \emph{per module}
through a pair of linear maps that re-express the adapter in $M_{v+1}$'s
activation geometry. Let
$P^{\mathrm{in}}_\ell \in \mathbb{R}^{d^{\mathrm{in}}_\ell \times d^{\mathrm{in}}_\ell}$
map an input activation in $M_{v+1}$'s space to the corresponding point in
$M_v$'s space, and
$P^{\mathrm{out}}_\ell \in \mathbb{R}^{d^{\mathrm{out}}_\ell \times d^{\mathrm{out}}_\ell}$
map an output activation in $M_v$'s space to $M_{v+1}$'s space. The transplanted
update is
\begin{equation}
  \Delta W'_\ell
  \;=\; P^{\mathrm{out}}_\ell \, \Delta W_\ell \, P^{\mathrm{in}}_\ell
  \;=\; \tfrac{\alpha}{r}
        \underbrace{\left(P^{\mathrm{out}}_\ell B_\ell\right)}_{B'_\ell}
        \underbrace{\left(A_\ell P^{\mathrm{in}}_\ell\right)}_{A'_\ell},
\label{eq:transplant}
\end{equation}
so $T(A)_\ell = (B'_\ell, A'_\ell)$ \emph{preserves the rank-$r$ LoRA
structure}: portability costs no extra inference parameters. Intuitively,
$A'_\ell$ first pulls an $M_{v+1}$ input back into $M_v$'s basis, the original
adapter operates there, and $P^{\mathrm{out}}_\ell$ pushes the result forward
into $M_{v+1}$'s output basis.

\paragraph{Fitting the maps.}
The maps are estimated from anchor activations only. Passing
$\mathcal{D}_{\mathrm{anchor}}$ through both checkpoints yields paired
input-activation matrices $H^{v}_\ell, H^{v+1}_\ell$ (and analogously output
matrices). We solve a ridge-regularized least-squares problem
\begin{equation}
  P^{\mathrm{in}}_\ell
  \;=\;
  \arg\min_{P}\;
  \big\lVert H^{v}_\ell - H^{v+1}_\ell P^\top \big\rVert_F^2
  + \lambda \lVert P \rVert_F^2 ,
\label{eq:procrustes}
\end{equation}
optionally constraining $P$ to be orthogonal (the orthogonal Procrustes
solution $P = UV^\top$ from the SVD $H_v^\top H_{v+1} = U\Sigma V^\top$), which
encodes the prior that version updates are approximately rotations of a shared
feature space. $P^{\mathrm{out}}_\ell$ is fit symmetrically with the roles of
$v$ and $v{+}1$ swapped. Both are closed-form, require a single forward pass of
$\mathcal{D}_{\mathrm{anchor}}$ through each base model, and never touch
$\mathcal{D}_{\mathrm{task}}$.

\subsection{Assumptions and Scope}
\label{sec:assumptions}
We assume (i) $M_v$ and $M_{v+1}$ share architecture and adapted-module
dimensions (the common case for minor version bumps within a family), and
(ii) a shared or alignable tokenizer so that $\mathcal{D}_{\mathrm{anchor}}$
induces token-aligned activations. We treat differing tokenizers, differing
hidden dimensions, and cross-family transfer as out of scope for this paper and
return to them in \S\ref{sec:limitations}.


\section{Method: Linear Representational Alignment}
\label{sec:method}
% Restate the transplant from the formulation; algorithm box: (1) run anchors
% through both bases, cache per-module in/out activations; (2) solve ridge /
% Procrustes for P^in, P^out per module; (3) apply Eq. (transplant); (4) load
% T(A) onto M_{v+1}. Variants: ridge vs orthogonal; shared-vs-per-layer maps;
% subset of modules; optional small nonlinear (MLP) map ablation.

\section{Representational Drift Across Versions}
\label{sec:drift}
% Empirical study: per-module CKA drift between M_v and M_{v+1} over the anchor
% set; drift vs depth; drift vs naive-transfer degradation (correlation plot).

\section{Experiments}
\label{sec:experiments}
% Setup: model pair(s), tasks, anchor set, metrics. Main table: naive | linear
% (ridge) | linear (Procrustes) | oracle, per task, with % gap recovered and
% compute. Ablations live in their own section.

\section{Ablations}
\label{sec:ablations}
% Anchor set size; ridge lambda; orthogonal vs unconstrained; which modules to
% align; per-layer vs shared map; linear vs nonlinear map; sensitivity to drift
% magnitude.

\section{Limitations}
\label{sec:limitations}
% Differing tokenizers / hidden dims / cross-family transfer out of scope
% (\S\ref{sec:assumptions}); single family pair; generative-quality eval caveats.

\section{Reproducibility}
\label{sec:repro}
% Code, anchor sets, fitted maps released. Lincoln Lab public-release clearance.

\bibliographystyle{plainnat}
\bibliography{references}

\end{document}
